\section{General-type models over a logarithmic base}\label{num:comparisons}

We now compare the variation of a log canonical model with sections on
the original total space. This comparison has two later uses. Marked
models detect which relative Iitaka bases are constant, while
Section~\ref{num:second-base} uses complete model systems to retain the
dimension of a second Iitaka base. In both applications, a finite base
change is only an auxiliary step: the sections must be tested on the
finite normalization of the original source.

Throughout this section, $k$ is an algebraically closed field of
characteristic zero and all varieties are projective. An effective
rational boundary has coefficients in $[0,1]$.
For a dominant morphism $p:H\to I$ between smooth integral varieties,
put $r=\dim H-\dim I$ and let $\mathcal V_p$ be the saturated kernel
of $dp$. Its canonical sheaf is
$(\bigwedge^r\mathcal V_p^\vee)^{**}$; at the generic point this is
the one-dimensional $k(H)$-space
$\bigwedge^r\Omega^1_{k(H)/k(I)}$. We compare the rank-one lattices
supplied by this sheaf and the ordinary relative canonical sheaf at
prime divisors of $H$. If $R_p$ is the divisorial zero of the pulled-back top
base differential, then taking determinant lattices gives
\begin{equation}\label{num:saturation}
 K_{\mathcal V_p}=K_{H/I}-R_p.
\end{equation}
Here $K_{H/I}=K_H-p^*K_I$, whereas $K_{\mathcal V_p}$ is the
canonical divisor of the saturated relative tangent sheaf. The term
$R_p$ records the difference between these two determinant lattices.
The saturated lattice of rational relative top forms depends only on
$k(I)\subset k(H)$ and the source model. It is therefore unchanged by
a birational change of the base. For a rational boundary $B$ on $H$,
write $B=C+B_v$, where $C$ is horizontal and $B_v$ is vertical over $I$.

\begin{lemma}[The inverse boundary]\label{num:boundary}
Let $D_I$ be a reduced SNC divisor on $I$. Then
\[
 R_p+(p^*D_I)_{\red}\geq p^*D_I.
\]
If $B$ is an effective rational boundary satisfying
$B\geq(p^*D_I)_{\red}$, it follows that
\begin{equation}\label{num:decomposition}
 K_H+B=K_{\mathcal V_p}+C+p^*(K_I+D_I)+V,
 \qquad V:=R_p+B_v-p^*D_I\geq0.
\end{equation}
\end{lemma}

\begin{proof}
Test a prime $E$ of $H$, with uniformizer $x$. In local boundary
coordinates on $I$, write $p^*t_j=x^{a_j}u_j$, with $u_j$ a unit.
Each logarithmic differential is $a_j\,dx/x+du_j/u_j$, and a
nonzero wedge contains at most one factor $dx/x$. If
$a=\operatorname{ord}_E(p^*D_I)>0$, the pulled-back ordinary base
volume thus vanishes to order at least $a-1$. This proves the first
inequality, also when $p(E)$ has higher codimension. If $a=0$, use
$R_p\geq0$. The second assertion follows from
Equation~\eqref{num:saturation} and the coefficient-one boundary
along the reduced inverse image of $D_I$.
\end{proof}

The boundary inequality supplies a fixed effective contribution from
the base. To use relative positivity for the remaining divisor, we
first arrange that all original vertical divisors have multiplicity
one after an alteration. Divisors introduced by a resolution are
allowed to remain exceptional: normality will remove them on the
finite normalization of the original source.

\begin{lemma}[Comparison at the original divisors]\label{num:alteration}
Let $C$ be an effective rational horizontal boundary on $H$, log smooth
over the generic point of $I$, and fix a finite extension of $k(I)$.
After enlarging it to a finite Galois extension $K'$, there are a finite
normalization $\tau:H^a\to H$ in a chosen compositum $k(H)K'$, a
smooth projective birational model $\mu:H'\to H^a$, and a morphism
$p':H'\to I'$ to a smooth projective model of $K'$. Put $\rho=\tau\mu$
and let $C'$ be the strict transform of the horizontal pullback of $C$.
They can be chosen so that:
\begin{enumerate}[label=\textup{(\roman*)}]
\item $C'$ has SNC support and the geometric generic pair is unchanged,
up to the field extension and choice of component;
\item every prime of $H'$ mapping onto a prime of the original $H$
is horizontal, or lies with multiplicity one over a prime of $I'$;
\item at every such prime, the relative-form lattices of
$K_{H'/I'}+C'$ and $\rho^*(K_{\mathcal V_p}+C)$ agree.
\end{enumerate}
Consequently, in every divisible degree $m$, a rational section of
$m(K_{H'/I'}+C')$ regular at these primes is an actual section of
$m\tau^*(K_{\mathcal V_p}+C)$ on $H^a$. The same assertion holds after adding
the pullback of a rational divisor on $H$. Every divisor above a
subset of codimension at least two in $I'$ is exceptional over $H$.
\end{lemma}

\begin{proof}
There are only finitely many prime divisors of $H$ outside a fixed
smooth log-family open. They include every multiple fiber component
and every divisor whose image on $I$ has codimension at least two.
Their restricted valuations can all be extracted on one smooth base
model. Indeed, if a divisorial valuation $v$ of $k(H)$ restricts
nontrivially to $w$ on $k(I)$, the residue-field inequality gives
\[
 \operatorname{trdeg}_k\kappa(w)
 \geq (\dim H-1)-r=\dim I-1.
\]
The valuation inequality gives the reverse bound. Thus $w$ is a
positive multiple of a divisorial valuation. Extract the finitely many
such valuations and resolve their common base model. Every original
vertical prime now has a divisorial center, which remains divisorial
on higher smooth base models.

All multiplicities in the next step are measured over this prepared
base. Normalize it in the prescribed field extension. For
each valuation above the selected base primes, take a root of a
uniformizer of degree equal to the least common multiple of the
relevant original multiplicities. Do this for all the finitely many
valuations and take a Galois closure. The ramification indices in these
towers multiply, so the final base index is divisible by every required
multiplicity; compare \cite[Lemmas~9.10--9.11]{KovacsPatakfalvi}.

The local calculation is at discrete valuation rings. If $t=ux^e$,
where $u$ is a unit, base ramification of index $n$ changes the
normalized fiber multiplicity to $e/\gcd(e,n)$. After an \'{e}tale
unit-root extension this is the normalization of $s^n=x^e$; separable
residue extensions do not change the indices. Thus $e\mid n$ makes
the multiplicity one, and further base ramification preserves it.
This also protects every original vertical prime outside the chosen
list, whose multiplicity was already one. Resolve the altered base,
the induced source map, and the boundary without changing the generic
smooth pair. The original divisorial centers and their DVRs retain the
properties just established.

At a vertical prime, choose separating relative residue parameters
$z_1,\ldots,z_r$. For the extracted base model, the ordinary relative
lattice has generator
$x^{-(e-1)}dz_1\wedge\cdots\wedge dz_r$ up to a unit: complete these
with base residue parameters and use
$dt=x^{e-1}(eu\,dx+x\,du)$. Saturation removes the factor
$x^{-(e-1)}$. The same residue differentials remain separating after
finite extension, and the new multiplicity is one. Hence the ordinary
relative lattice upstairs is the pullback of the original saturated
lattice. At a horizontal prime the DVR contains $k(I)$ as a field,
so finite separable base extension is \'{e}tale there. Its relative
lattice and boundary coefficient are unchanged. This proves (i)--(iii).

Every prime of $H^a$ lies over a prime of $H$, since $\tau$ is finite.
Its strict transform is therefore among the primes just tested.
Normality of $H^a$ gives the asserted global section. The same tests
hold after tensoring by a pulled-back line bundle. Property (ii) proves
the final assertion. Poles at the remaining resolution-exceptional
primes are harmless because the section is asserted on $H^a$, not on
that resolution.
\end{proof}

The finite-generation and log-canonical-model statements from
\cite{BCHM}, whose cited version is formulated over $\C$, are used
over $k$ by the following field comparison. Descend the projective
morphism, the rational boundary, a log resolution, and an ample
decomposition witnessing relative bigness to a field $k_0$ finitely
generated over $\Q$. Enlarge $k_0$ to define all this finite data.
An embedding $k_0\hookrightarrow\C$ preserves the geometric klt
and bigness conditions, so the complex theorem applies. Formation
of each graded direct-image sheaf in the relative adjoint ring
commutes with field extension. Finite generation descends through
the faithfully flat extension $k_0\subset\C$: finitely many
homogeneous generators upstairs involve finitely many elements of
the original graded algebra, and faithful flatness detects the
vanishing of the resulting cokernel. Applying this on an affine
cover of the target gives finite generation over $k_0$, hence over
$k$. Relative $\operatorname{Proj}$ commutes with field extension,
so the resulting canonical model and its comparison maps have the
same compatibility. We use this argument also over the parameter
fields in Section~\ref{sec:marking}.

We can now apply positivity for families of general-type pairs
without changing the section problem on $H$. The variation below is
the variation of the log canonical \emph{model together with its
boundary}; it is a different invariant from birational descent of the
whole geometric generic fiber.

\begin{proposition}[Model-pair variation over a logarithmic base]
\label{num:model-pair}
Let $B$ be an effective rational SNC boundary on $H$, and let $D_I$
be a reduced SNC divisor on $I$, such that
\[
 B\geq(p^*D_I)_{\red},\qquad \kappa(I,K_I+D_I)\geq0.
\]
Suppose the log canonical divisor of a geometric generic fiber
component is big. Then $\kappa(H,K_H+B)\geq r$. If that generic pair
is klt, then
\begin{equation}\label{num:KP}
 \kappa(H,K_H+B)\geq r+\operatorname{Var}(p,C),
 \qquad C=B_{\mathrm{hor}},
\end{equation}
Here $\operatorname{Var}(p,C)$ is the minimum transcendence degree
over $k$ of an algebraically closed field over which the geometric
generic log canonical \emph{model pair}, including its boundary, is
isomorphic to a base extension. Thus it measures descent of that
marked model; the invariant $\Var(f)$ in
Equation~\eqref{eq:variation-definition} measures birational descent
of the whole geometric generic fiber.
In the klt case, if the left side is at most $r$, this model pair
becomes constant after a finite extension of the base function field. For disconnected
geometric generic fibers, the model and its variation refer to a chosen
component after adjoining the Stein field.
\end{proposition}

\begin{proof}
If $r=0$, the chosen geometric generic component is a point and its
variation is zero. The saturated relative tangent sheaf has rank zero,
so $K_{\mathcal V_p}=0$, and there is no horizontal boundary.
Equation~\eqref{num:decomposition} and a nonzero plurisection of
$K_I+D_I$ give the required nonnegative Iitaka dimension directly.
The model is a point after the Stein extension. We may therefore
assume $r>0$.

First suppose the generic pair is klt. Apply
Lemma~\ref{num:alteration}, prescribing the algebraic closure of
$k(I)$ in $k(H)$ and choosing the corresponding component. The altered
geometric generic fiber is integral and has the same klt big pair.
The relative form of Kov\'acs--Patakfalvi's inequality,
\cite[Theorem~9.9]{KovacsPatakfalvi} with its auxiliary base divisor
$M=0$, gives
\[
 \kappa(H',K_{H'/I'}+C')\geq r+\operatorname{Var}(p,C).
\]
Its hypotheses are satisfied by the smooth projective SNC pair and
its klt big generic log divisor; no condition on $\kappa(I')$ is
required. The model-pair convention is
\cite[Definition~6.16, Remark~6.17, Corollary~6.20,
Definition~9.3 and Remark~9.4]{KovacsPatakfalvi}; the
canonical model exists for klt big pairs by finite generation
\cite{BCHM}. Finite extension does not change this variation.

By Lemma~\ref{num:alteration}, these systems give sections of
$m\tau^*(K_{\mathcal V_p}+C)$ on the finite normalization of the
original $H$. Their ratios are unchanged. Lemma~\ref{num:finite}
therefore gives the same lower bound for $K_{\mathcal V_p}+C$ on $H$.
Choose $0\ne\alpha\in H^0(I,\ell(K_I+D_I))$. In common divisible
degrees $m$, multiply by $p^*\alpha^{m/\ell}$ and the canonical
section of $mV$ in Equation~\eqref{num:decomposition}. These common
factors embed the systems into those of $m(K_H+B)$ without changing
their rational-map dimensions. This proves Equation~\eqref{num:KP}.

For a generic lc pair, decrease the horizontal coefficients by a
sufficiently small rational factor. Its log divisor stays big, while
the generic pair becomes klt. The required vertical boundary is
unchanged. The resulting systems are contained in those for $B$, which
proves the lower bound $r$; no assertion about the variation of the
unperturbed lc pair is needed.

In the zero-excess klt case, the inequality forces model-pair variation
zero. A stable general-type pair has finite automorphism group. The
constant moduli point has a representative over $k$, and its generic
isomorphism torsor has a point after a finite field extension. This
makes the model pair constant, with its boundary retained.
\end{proof}
